\SourceLeaf{217}%
\bigskip\Needspace{12\baselineskip}
\pdfbookmark[1]{Further Report on a New Planet Discovered by Dr. Olbers. June 1802}{pallas180206}
\fancyhead[L]{\small PALLAS}
\begin{center}
{\small FURTHER REPORT ON A NEW AND MOST REMARKABLE PLANET OF OUR SOLAR SYSTEM, DISCOVERED BY DR. OLBERS IN BREMEN AND MOVING BETWEEN MARS AND JUPITER.\par}
\medskip\hrule\medskip
Monatliche Correspondenz. Vol. V. pp.\ 591—606. June 1802.\par
\medskip\hrule
\end{center}\medskip
When we first reported the remarkable discovery of \textsc{Olbers}'s new celestial body in the preceding issue, the observations, calculations and experience gathered up to that point did not yet entitle us to announce this strange wandering star as a planet belonging to our solar system. We had, however, already voiced the conjecture that this body, so distinctive in its motion and appearance, did not belong to the class of ordinary so-called comets. Yet the suspicion of a \textit{planetary nature} was already so strong at the very first sight of this celestial body that its discoverer ventured at once to try to fit its orbit to a circle concentric with the Sun, without first trying a parabola as is usually done. But as we already reported on p.\ 486 of the previous issue, no circular orbit could represent the available observations at all. Few as they were, the body's large motion in latitude already made it possible to rule out the circle for good in future investigations of its orbit.

Dr.\ \textsc{Gauss} tried the same thing. He too confirmed Dr.\ \textsc{Olbers}'s finding that \textit{Pallas} was always moving faster than \textsc{Kepler}'s laws would allow it to move in a circle. But as soon as Dr.\ \textsc{Gauss} received the observation of this body that Dr.\ \textsc{Olbers} had made on 17 April, he undertook to determine the conic section independently of any hypothesis, using his own method, whose excellence \textit{Ceres} had already confirmed so brilliantly. It was foreseeable that the observations made with the ring micrometer would lack the accuracy and duration needed for this delicate calculation, especially since \textit{Pallas} was in an unfavorable position at the beginning of April, with the tangent to its geocentric motion passing through the position of the Sun. Nevertheless, Dr.\ \textsc{Gauss} found an ellipse whose main dimensions differed little from those of the one determined later. The result was in broad agreement with the later one, showing that Pallas was a planet between Mars and Jupiter whose orbit at one point comes very close to the orbit of Ceres and has a considerable eccentricity as well as a very large inclination.

The eccentricity in this case was $0.3$, greater than that of all the other planets, yet too small to justify denying \textit{Pallas} the name of a planet. The ratio of the major to the minor axis is $1:\BareRoot{(1-(0.3)^2)}$, or approximately $1:0.95$, so the ellipse is still less flattened than Jupiter's disk.

Since the observations used for this purpose are naturally far less accurate than meridian observations, Dr.\ \textsc{Gauss} had good grounds to distrust these elements and withheld their publication until they could be checked further. Only after he received our \textit{Seeberg} meridian observations of this body on 15, 18 and 19 April, and later three Paris observations on 10, 12 and 13 April, did he decide, in the absence of a longer series of accurate observations, to improve these prelimi%
\SourceLeaf{218}%
nary elements so that they would fit six \textit{Seeberg} observations on 4, 5, 7, 15, 18 and 19 April as well as the three \textit{Paris} observations. In this way he found the following at once on his first attempt
\expandafter\def\csname HistoricalNoteText19\endcsname{Dr.\ \textsc{Gauss} used exactly the same assumption for \textit{Ceres}. He treated the aphelion and $\AscendingNode$ as fixed relative to the fixed stars rather than to our equinox.}
\Needspace{20\baselineskip}
\begin{center}
\textit{Elements of Pallas} (I).\par\medskip
[Handwriting:] Heliocentric arc $4^0\ 11'$\par\medskip
Orbit in a prograde ellipse.\par\medskip
{\small\setlength{\tabcolsep}{2pt}
\begin{tabular}{l r}
Epoch 1802 March 31, noon at Seeberg $\ .\ .\ .\ .\ .\ .\ .$&$166^0\quad1'\ 37''2$\\
Mean daily tropical motion $\ .\ .\ .\ .\ .\ .\ .$&$800''770$\\
Orbital period $\ .\ .\ .\ .\ .\ .\ .\ .\ .\ .\ .\ .\ .\ .\ .\ .\ .$&$1618\tfrac12$ days\\
Logarithm of the semimajor axis $\ .\ .\ .\ .\ .\ .\ .$&$0.4310494$
\end{tabular}
\par\smallskip
$\left.\begin{array}{l}\text{Aphelion}\\\text{Node}\end{array}\right\}\quad
\text{\shortstack{at the epoch\\fixed relative to the stars\footnotemark}}\left\{\begin{array}{lr}.&304^0\ 36'\ 30''\\.&172^0\quad9'\ 58''\end{array}\right.$
\footnotetext{\csname HistoricalNoteText19\endcsname}
\par\smallskip
\begin{tabular}{l r}
Eccentricity $\ .\ .\ .\ .\ .\ .\ .\ .\ .\ .\ .\ .\ .\ .\ .\ .$&$0.215708$\\
Inclination $\ .\ .\ .\ .\ .\ .\ .\ .\ .\ .\ .\ .\ .\ .\ .\ .\ .$&$33^0\ 39'\ 16''6$
\end{tabular}}
\end{center}
Dr.\ \textsc{Gauss} now compared these elements very closely with the \textit{Seeberg} observations he had at the time and those he received from us later, as well as with the \textit{Paris} observations and one by \textsc{Olbers}. The table below shows how closely they agreed.
\begin{center}
{\small\setlength{\tabcolsep}{3pt}
\begin{tabular}{r r|r|r|r|r|l}
\multicolumn{2}{c|}{}&&&\multicolumn{2}{c|}{Differences}&\\
\multicolumn{2}{c|}{1802}&\shortstack{Calculated\\R.A. of Pallas}&\shortstack{Calculated\\declination\\of Pallas}&in R.A.&in decl.&Observer\\\hline
April&4&$183^0\ 44'\quad6''3$&$13^0\ 54'\ 51''0$&$-0''3$&$-2''0$&von Zach\\
&5&$183\ 34\ 25.7$&$14\ 13\ 17.2$&$+2.0$&$-5.7$&—\\
&7&$183\ 15\ 40.0$&$14\ 48\ 55.5$&$+1.5$&$-6.6$&—\\
&8&$183\quad6\ 35.2$&$15\quad6\quad9.1$&$-2.6$&$-0.9$&—\\
&10&$182\ 48\ 56.5$&$15\ 39\ 40.9$&$-1.5$&$-1.1$&Dr.\ Burckhardt\\
&12&$182\ 32\ 34.3$&$16\ 11\quad4.2$&$+2.8$&$+6.2$&Méchain\\
&13&$182\ 24\ 47.2$&$16\ 26\quad6.6$&$+2.8$&$+0.6$&—\\
&15&$182\ 10\ 13.2$&$16\ 54\ 33.0$&$-3.3$&$+2.2$&von Zach\\
&18&$181\ 50\ 26.2$&$17\ 34\ 26.5$&$-4.4$&$\ .\ .\ .$&—\\
&19&$181\ 44\ 28.0$&$17\ 46\ 51.7$&$+2.7$&$-2.5$&—\\
&24&$181\ 19\ 39.4$&$18\ 42\ 27.8$&$+1.6$&$\ .\ .\ .$&—\\
&25&$181\ 15\ 44.1$&$18\ 52\ 18.3$&$+11.9$&$\ .\ .\ .$&—\\
&26&$181\ 12\ 10.4$&$19\quad1\ 43.5$&$+8.6$&$-5.9$&—\\
&27&$181\quad8\ 57.9$&$19\ 10\ 44.0$&$+12.2$&$-2.8$&—\\
&29&$181\quad3\ 39.6$&$19\ 27\ 31.2$&$+23.0$&$-12.5$&—\\
&30&$181\quad1\ 33.9$&$19\ 35\ 18.2$&$+24.0$&$-15.3$&—\\
May&1&$180\ 59\ 48.1$&$19\ 42\ 42.5$&$+25.6$&$+1.7$&—\\
&1&$180\ 59\ 38.0$&$19\ 43\ 38.0$&$+20.0$&$+7.0$&Dr.\ Olbers
\end{tabular}}
\end{center}
Adding the differences with their signs reversed to the calculated values gives the observations.

Here, as with \textit{Ceres}, Dr.\ \textsc{Gauss} keeps to his old practice, which deserves to be followed. He derives the right ascensions and declinations from the elements themselves and compares them with the observed ones, without using the observed R.A. and declination \textit{to derive observed longitudes and latitudes}. The reason for this procedure is so evident that one has to wonder why
\SourceLeaf{219}%
it has not long since been universally adopted and followed. The effort is the same whether one takes the observation a step toward the calculation or carries the calculation all the way to the observation. Using a procedure of Dr.\ \textsc{Gauss}'s own, this parallax calculation in fact becomes somewhat simpler. The chief advantage, however, is that one can separate the good observations from those that may be less reliable. Otherwise an inaccurate declination can easily spoil a good right ascension and give both an erroneous latitude and a poor longitude. Since the differences in the above comparison of the observations with the orbital elements keep increasing toward the end, one could foresee that, given the short span of the observations, a difference of $25''$ must still have a very strong effect on the determination of these elements. The further investigation that Dr.\ \textsc{Gauss} undertook immediately after this comparison had also shown him that the 15 days of observations on which his previous elements were based were insufficient to determine the orbit accurately at that point, though they already indicated the nature of the body fairly clearly. He then immediately determined new elements from the whole series of our \textit{Seeberg} observations, fitting them as closely as possible to those observations from 4 April to 1 May. As the arc already traversed along the orbit amounts to $7^0\ 26'$, at least as far as its planetary nature is concerned, there is probably no longer any reason for doubt about \textsc{Olbers}'s celestial body. The new ellipse of this planetary body is accordingly as follows.
\begin{center}
\textit{Elements of Pallas from 27 days of Seeberg observations.}\par\medskip
(II.)\par\medskip
{\small\setlength{\tabcolsep}{2pt}
\begin{tabular}{l r}
Epoch 1802 March 31, noon at Seeberg $\ .\ .\ .\ .\ .\ .\ .$&$161^0\ 12'\ 43''2$\\
Mean daily tropical motion $\ .\ .\ .\ .\ .\ .\ .\ .$&$757''166$\\
Logarithm of the semimajor axis $\ .\ .\ .\ .\ .\ .\ .\ .$&$0.4472636$
\end{tabular}
\par\smallskip
$\left.\begin{array}{l}\text{Aphelion}\\\text{Node}\end{array}\right\}\ \text{fixed relative to the stars}\left\{\begin{array}{lr}
.\ .\ .\ .\ .\ .\ .\ .\ .&300^0\quad5'\quad4''\\
.\ .\ .\ .\ .\ .\ .\ .\ .&172\ 34\ 35
\end{array}\right.$
\par\smallskip
\begin{tabular}{l r}
Eccentricity $\ .\ .\ .\ .\ .\ .\ .\ .\ .\ .\ .\ .\ .\ .\ .\ .$&$0.2591096$\\
Inclination $\ .\ .\ .\ .\ .\ .\ .\ .\ .\ .\ .\ .\ .\ .\ .\ .\ .\ .$&$35^0\quad0'\ 42''$
\end{tabular}}
\end{center}

\Needspace{18\baselineskip}
Dr.\ \textsc{Gauss} then compared the whole series of our \textit{Seeberg} observations with these elements and found the agreement shown below.
\begin{center}
{\small\setlength{\tabcolsep}{5pt}
\begin{tabular}{r r|r|r|r|r}
\multicolumn{2}{c|}{}&&&\multicolumn{2}{c}{Differences}\\
\multicolumn{2}{c|}{1802}&\shortstack{Calculated\\R.A.\\of Pallas}&\shortstack{Calculated north\\declination\\of Pallas}&in R.A.&in decl.\\\hline
April&4&$183^0\ 44'\quad5''7$&$13^0\ 54'\ 55''4$&$-0''9$&$+3''4$\\
&5&$183\ 34\ 24.6$&$14\ 13\ 22.6$&$+0.9$&$-0.3$\\
&7&$183\ 15\ 39.6$&$14\ 49\quad1.3$&$+1.6$&$-0.7$\\
&8&$183\quad6\ 35.4$&$15\quad6\ 15.2$&$-2.4$&$+5.2$\\
&15&$182\ 10\ 17.9$&$16\ 54\ 35.7$&$+1.4$&$+4.9$\\
&18&$181\ 50\ 29.4$&$17\ 34\ 27.1$&$-1.2$&$\ .\ .\ .$\\
&19&$181\ 44\ 30.7$&$17\ 46\ 51.9$&$+5.4$&$-2.5$\\
&24&$181\ 19\ 35.8$&$18\ 42\ 27.9$&$-2.0$&$\ .\ .\ .$\\
&25&$181\ 15\ 38.3$&$18\ 52\ 18.9$&$+6.1$&$\ .\ .\ .$\\
&26&$181\ 12\quad2.0$&$19\quad1\ 44.9$&$+0.2$&$-4.5$\\
&27&$181\quad8\ 46.6$&$19\ 10\ 46.3$&$+1.0$&$-0.5$
\end{tabular}
\par\medskip
\begin{tabular}{r r|r|r|r|r}
\multicolumn{2}{c|}{}&&&\multicolumn{2}{c}{Differences}\\
\multicolumn{2}{c|}{1802}&\shortstack{Calculated\\R.A.\\of Pallas}&\shortstack{Calculated north\\declination\\of Pallas}&in R.A.&in decl.\\\hline
April&29&$181^0\quad3'\ 21''2$&$19^0\ 27'\ 36''2$&$+4''6$&$-7''5$\\
&30&$181\quad1\ 11.4$&$19\ 35\ 24.9$&$+1.5$&$-8.6$\\
May&1&$180\ 59\ 22.9$&$19\ 42\ 50.8$&$+0.4$&$+10.0$\\
&2&$180\ 57\ 56.7$&$19\ 49\ 53.2$&$+4.8$&$+1.6$\\
&3&$180\ 56\ 52.4$&$19\ 56\ 32.9$&$-3.2$&$-14.8$\\
&5&$180\ 55\ 49.6$&$20\quad8\ 45.6$&$-5.0$&$+6.9$\\
&6&$180\ 55\ 51.4$&$20\ 14\ 19.4$&$-2.7$&$\ .\ .\ .$\\
&7&$180\ 56\ 14.9$&$20\ 19\ 32.1$&$-8.9$&$-3.2$\\
&8&$180\ 57\quad0.4$&$20\ 24\ 23.7$&$-9.6$&$+3.6$\\
&11&$181\quad1\ 25.5$&$20\ 36\ 58.8$&$-16.2$&$-14.4$
\end{tabular}}
\end{center}
\SourceLeaf{220}%

Dr.\ \textsc{Gauss} could very easily have obtained a slightly closer fit to the last right ascensions. But because the declinations were less reliable as twilight drew near, he could not have been sure that this would really yield an \textit{improvement} in the orbit. He therefore intends to wait for later observations before correcting these elements again.

Be that as it may, all the observations, calculations and comparisons so far seem to establish the planetary nature of \textit{Pallas} and to protect us against all doubts, misgivings and indefinite hypotheses. As Dr.\ \textsc{Gauss} aptly remarks in a letter, “A precise calculation can be countered only by an equally precise calculation. Empty conjectures and vague reasoning cannot answer it. To find Pallas again easily in 1803, my elements II will of course still need considerable improvements, which they will hopefully receive if only enough further observations can be made. But these improvements will be \textit{considerable} only for determining the position in 1803. They have no such significance for the \textit{planetary} or \textit{nonplanetary nature} of \textit{Pallas}.”

Even with the earlier elements, the orbit of \textit{Pallas} came very close to that of Ceres where its ascending node lies on the Ceres orbit. This led Dr.\ \textsc{Olbers} to conjecture that the orbits of \textit{Pallas} and \textit{Ceres} actually intersect, perhaps looping into one another like two links of a chain. From a rough calculation of the distances of \textit{Pallas} and \textit{Ceres} along the line of nodes of their orbits, Dr.\ \textsc{Gauss} found that at $\DescendingNode$ the distance from the Sun was $2.86$ for \textit{Pallas} and $2.93$ for \textit{Ceres}. At $\AscendingNode$ the distances are less nearly equal. But the series of observations is still too short to rule out the possibility that the distances at this point may be exactly equal. Dr.\ \textsc{Gauss} intends to let the calculation take its full course without introducing anything hypothetical at all. It is of course too early to dwell on hypotheses, since if we manage to find \textit{Pallas} again next year, we will soon have certainty in their place.

Meanwhile, what speculations about the origin and history of our planetary system will this strange celestial body give rise to? What an unexpectedly large inclination for a planet! And how striking the position of its orbit relative to the orbit of Ceres! This suggested to Dr.\ \textsc{Olbers} that the two planets might perhaps be mere \textit{remnants} of a single planet shattered by the impact of a comet. Who would ever have expected such a thing in our planetary system? And what new, important and far-reaching insights can and will this celestial body lead us to? In more than one respect it ranks among the most important astronomical discoveries ever made.

Dr.\ \textsc{Gauss} finds the close agreement between the mean motions of these two planets just as remarkable. So far it is impossible to decide with certainty whether they might be \textit{exactly} equal. After all, we still do not know the actual daily. motion of \textit{Ceres} to within one second. Yet small as the masses of \textit{Ceres} and \textit{Pallas} are, there must be a limit within which the difference between their mean motions cannot fall without being reduced entirely to zero by their mutual action. According to elements II, these mean motions differ only slightly, and perhaps a more accurate knowledge of the orbits will bring them much closer together. It might therefore be that \textit{Ceres} and \textit{Pallas} have the same orbital period and are never very far from each other, yet have never been too close since the beginning and cannot come too close in the future. Because their equation of the center is so different, at $\AscendingNode$ \textit{Pallas} would always be ahead, and at $\DescendingNode$ \textit{Ceres} would always be ahead. After a long series of years one will then also be able to determine very well, from observations of \textit{Ceres}, the
\SourceLeaf{221}%
mass of \textit{Pallas}, and from observations of \textit{Pallas}, the mass of \textit{Ceres}. Dr.\ \textsc{Gauss} finds from a rough calculation that a difference of 12 seconds in the mean daily motion gives a synodic revolution of nearly three hundred years.

\begingroup\rightskip=0pt plus1em\relax
What investigations will such a unique phenomenon prompt? What new needs will it awaken in science? For all we knew before, we would never have expected such extraordinary phenomena in the system of the world! How our customary methods of calculation will fail us and prove inadequate! For example, consider even just the calculation of the perturbations of this new celestial body. In the only methods in use up to now, the eccentricity and inclination of orbits have been treated as if they were infinitesimal quantities. The expansion was made in their powers, neglecting the higher ones. This will no longer be possible for two such large elements of the Pallas orbit, whose higher powers certainly cannot be neglected. Perhaps \textit{Pallas} therefore necessarily opens a new epoch in physical astronomy as well. What we know by rigorous calculation is, after all, the only thing we \textit{really know} about the structure of the world. But the person doing the calculations needs experience and observations. Only through them can one stamp hypotheses as truth. Hence all astronomers have observed \textit{Pallas} with as much diligence as care. We give our readers the whole collected series of these precious observations. Our journal has so far been the archive of all observations of \textit{Ceres}. It is to serve the same purpose for \textit{Pallas}, and we will spare no effort to make this collection as complete as possible, so that astronomers may find everything together here that they would otherwise have to hunt for in scattered publications, and would find only too late, or perhaps never.\par\endgroup

—\quad —\quad —

To make it easier for astronomers equipped with large equatorial instruments to find \textit{Pallas}, Dr.\ \textsc{Gauss} has used his orbital elements II to calculate the following ephemeris of its motion. It may seem somewhat bold to try to use elements based on only 27 days of observations to determine the position of the planet 2 months in advance. The ephemeris is sure to differ quite noticeably from the sky toward the end. But since the discrepancy is likely to be only a few minutes, it will still be useful for finding the planet easily, especially since one can already foresee that the error in R.A. (and presumably also in declination) will be negative.
\Needspace{12\baselineskip}
\begin{center}
\textit{Ephemeris for Pallas Olbersiana, for midnight at Seeberg.}\par\medskip
{\small\setlength{\tabcolsep}{5pt}
\begin{tabular}{r r|r|r@{\qquad}r r|r|r}
\multicolumn{2}{c|}{1802}&\shortstack{R.A.\\of Pallas}&\shortstack{Declination\\of Pallas}&\multicolumn{2}{c|}{1802}&\shortstack{R.A.\\of Pallas}&\shortstack{Declination\\of Pallas}\\\hline
May&24&$181^0\ 57'$&$21^0\quad1'$ N.&June&14&$185^0\ 12'$&$20^0\ 17'$ N.\\
&27&$182\ 18$&$21\quad0$&&17&$185\ 48$&$20\quad5$\\
&30&$182\ 41$&$20\ 57$&&20&$186\ 27$&$19\ 52$\\
June&2&$183\quad6$&$20\ 52$&&23&$187\quad7$&$19\ 37$\\
&5&$183\ 34$&$20\ 46$&&26&$187\ 49$&$19\ 22$\\
&8&$184\quad5$&$20\ 38$&&29&$188\ 32$&$19\quad6$\\
&11&$184\ 37$&$20\ 28$&&&&
\end{tabular}}
\end{center}
\begin{flushright}[\textsc{von Zach}.]\end{flushright}
